\section{Hyperparameter Tuning}
\label{sec:tuning}

\textbf{Every model gets the same search, and tuning never sees validation or test
data} (\repofile{src/int4096/experiments/tuning.py}{tuning code}).

\begin{itemize}
  \item \textbf{Search:} random search, \TrialsUsed{} trials and at most
  \SecondsUsed{}~s per run (protocol caps: \TrialsMax{} and \SecondsMax{}~s), seed
  42. Repeated draws count as trials.
  \item \textbf{Score:} mean error on the price change across expanding-window
  cross-validation \textbf{inside the training block}: five folds, a one-row gap,
  minimum initial fraction 0.30.
  \item \textbf{Final fit:} the best setting is refitted on the whole training
  block, then scored once on validation. Every attempt is
  \repodir{experiments/models}{logged}.
  \item \textbf{Early stopping:} gradient boosting uses a random 10\% of its
  training rows, so its iteration setting is an upper limit.
\end{itemize}

% GENERATED by scripts/build_report_tuning.py -- do not edit by hand.
\begin{table}[H]
\centering
\caption{Search grids, distinct settings tried and chosen values (standardised inputs, P2).}
\label{tab:tuning}
\footnotesize
\setlength{\tabcolsep}{3pt}
\renewcommand{\arraystretch}{1.15}
\begin{tabular}{@{}L{0.15\linewidth} L{0.34\linewidth} L{0.09\linewidth} L{0.18\linewidth} L{0.18\linewidth}@{}}
\toprule
\textbf{Model} & \textbf{Search grid} & \textbf{Tried} & \textbf{Chosen, AAPL} & \textbf{Chosen, BTC} \\
\midrule
Linear baseline & $\alpha$ \{0.1, 1, 10\} & 3 of 3 & $\alpha$ 10 & $\alpha$ 10 \\
Ridge & $\alpha$ \{0.01, 0.1, 1, 10, 100, 1000\} & 6 of 6 & $\alpha$ 1000 & $\alpha$ 1000 \\
SVR (RBF) & C \{0.1, 1, 10, 100, 1000\}; $\varepsilon$ \{0.001, 0.01, 0.1\}; $\gamma$ \{scale, 0.001, 0.01, 0.1\} & 10 of 60 & C 0.1, $\varepsilon$ 0.1, $\gamma$ 0.01 & C 10, $\varepsilon$ 0.01, $\gamma$ 0.1 \\
Random forest & trees \{100, 200\}; depth \{none, 4, 8, 16\}; leaf \{1, 5, 20\} & 9 of 24 & trees 100, depth none, leaf 20 & trees 100, depth 4, leaf 20 \\
Hist.\ gradient boosting & rate \{0.01, 0.05, 0.1\}; iterations \{100, 300\}; depth \{none, 3, 5\}; leaf \{5, 20, 50\} & 10 of 54 & rate 0.01, iterations 300, depth 3, leaf 20 & rate 0.01, iterations 300, depth 3, leaf 20 \\
\bottomrule
\end{tabular}
\end{table}


\paragraph{What tuning showed.}
\begin{itemize}
  \item \textbf{Coverage differs:} both linear grids were tried in full; the other
  three were only sampled.
  \item \textbf{Grid edge:} Ridge chose the \textbf{largest penalty} on both assets,
  so a wider grid is a sensible next step.
  \item \textbf{Simpler settings often won}, such as strong penalties, large forest
  leaves and a low boosting rate. This is consistent with weak signal but does not
  prove it.
  \item \textbf{Scaling and choice:} without scaling (P0), only BTC SVR selected a
  different setting.
\end{itemize}
